%--
\chapter{最小生成树}

%----------------------------------------
\section*{授课提纲}

\nid MST的定义：
%--
\begin{itemize}
    \item 更贴近人的定义 vs. 更贴近机器的定义（MST Property）
    \item 二者等价性的证明
\end{itemize}

\nid \alg{Prim}算法的设计与实现：
%--
\begin{itemize}
    \item \alg{Prim}算法设计 $\Rightarrow$ \alg{Prim}算法正确性证明（基于MST Property）
    \item \alg{Prim}算法实现（基于PQ）$\Rightarrow$ \alg{Prim}算法代价分析（基于抽象PQ的抽象代价分析）
\end{itemize}

\nid \alg{Kruskal}算法的设计与实现：
%--
\begin{itemize}
    \item \alg{Kruskal}算法设计 $\Rightarrow$ \alg{Kruskal}算法正确性证明
    \item \alg{Kruskal}算法实现 $\Rightarrow$ \alg{Kruskal}算法代价分析（基于抽象并查集的代价分析）
\end{itemize}

\nid \alg{Prim} vs. \alg{Kruskal}：
%--
\begin{itemize}
    \item 实例化抽象数据类型的实现
    \item 按$n$和$m$的关系，分情况讨论
\end{itemize}

\nid \alg{BestFS}：
%--
\begin{itemize}
    \item 实例化为\alg{Prim}
    \item 实例化为\alg{Dijkstra}
\end{itemize}

